\section{Notation and Linearization}\label{app:encoding}
Table~\ref{tab:notation} summarizes the objects used in the formulation. All finite ranges, time horizons, and coefficient versions are part of a model instance.

\begin{table}[htbp]
\centering\small
\begin{tabular}{p{.20\linewidth}p{.72\linewidth}}
\toprule
Symbol & Meaning \\
\midrule
$\cW,\cD_W$ & Required deployment scenarios and feasible bounded implementation domain \\
$R_k,\cF_k,\cR_k$ & Active representation, feasible assignments, and decoded realizable designs \\
$a,p,m,\sigma$ & Algorithmic organization, program, mapping, and execution organization \\
$h,b$ & Shared hardware structure and software-visible boundary choices \\
$q_{rc},x_{rct}$ & Implementation selection and time-indexed start decisions \\
$y_{ip},b_{j\mu}$ & Hardware-slot configuration and boundary-mode selection \\
$U_{rch},S_{rc}$ & Local resource-demand and internal-scratch profiles \\
$\ell_{v\tau},B_v$ & Value residency and fixed byte footprint \\
$L_w,\widehat J,J$ & Scenario latency, modeled objective, and target evaluator objective \\
\bottomrule
\end{tabular}
\caption{Main notation. A symbol denoting a program $p$ is distinct from the configuration index $p\in P_i$ in hardware-slot sums.}\label{tab:notation}
\end{table}

\textbf{Logical conjunction.}
For binary selections $q,y$, an auxiliary $z=q\land y$ is encoded by
\begin{equation}
 z\le q,\quad z\le y,\quad z\ge q+y-1,\quad z\ge0.
\end{equation}
This is useful for configuration-dependent duration or energy coefficients. It enumerates local interactions without requiring an opaque candidate for every whole-system combination.

\textbf{Binary--continuous products.}
For $z=qv$ with $q\in\{0,1\}$ and known bounds $L\le v\le U$, the exact linearization is
\begin{equation}
 Lq\le z\le Uq,\qquad v-U(1-q)\le z\le v-L(1-q).
\end{equation}
Tight bounds matter for relaxation strength. A generic large constant can make a valid model unnecessarily difficult to solve.

\textbf{Residency.}
Residency variables cannot be free capacity switches. For a selected implementation, they are constrained to remain active from production or preallocation until every required consumer completes. One finite encoding enumerates allowed start/last-use pairs and derives the active interval; another uses production and release events with cumulative conservation. Aliasing and buffer reuse require nonoverlapping live ranges or explicit ordering. Internal scratch and exported values must not be counted twice.

\textbf{Configuration-dependent timing.}
A single duration coefficient is valid only for the hardware conditions under which it was defined. A finite $(\text{implementation},\text{hardware mode})$ table can use conjunction variables. A factorized model instead represents the relevant service capacities explicitly and derives event times. The former is compact when interactions are local; the latter is needed when relative order and interference change with the configuration.

\textbf{Conditional families.}
Suppose an attention region has family selectors $g_{\mathrm{mat}},g_{\mathrm{stream}}$ summing to one. Each internal operation $u$ belonging to family $a$ satisfies $\sum_{r,c:\,u\in\operatorname{cover}_a(r,c)}q_{rc}=g_a$. The selected family exports the same region-level inputs and output. Transfers and lifetimes are internal to that family except for shared external values. This avoids treating an online recurrence as a cover of intermediate tensors that it intentionally does not materialize.

\section{A Complete Finite Boundary-Choice Model}\label{app:boundary}
Consider a region that must transfer and consume $N$ tiles in each workload scenario. Let $\mathcal H$ be a finite set of machine configurations and $\mathcal M=\{\mathrm{sw},\mathrm{async},\mathrm{ctrl}\}$ the permitted transfer-control modes. Each hardware/mode pair has a generator and compiler interface. Let $\mathcal C_w$ contain software variants with fixed tile sizes, traversal order, and buffering depth for scenario $w$.

The coefficients are: shared area $A_h$, incremental mode area $A^B_{h\mu}$, compatibility $K_{wh\mu c}\in\{0,1\}$, and evaluated or modeled latency $t_{wh\mu c}$. The mode-area coefficient includes only components not already counted in $A_h$. Compatibility includes numerical eligibility, supported traversal, live storage, resource demand, and interface support. Thus a specialized controller may support fewer variants than the programmable mode.

Choose binary hardware $y_h$, mode $b_\mu$, and software $q_{wc}$ variables with
\begin{equation}
 \sum_h y_h=1,\qquad \sum_\mu b_\mu=1,\qquad
 \sum_{c\in\mathcal C_w}q_{wc}=1\quad\forall w.
\end{equation}
Let $v_{h\mu}=y_h\land b_\mu$ and $z_{wh\mu c}=v_{h\mu}\land q_{wc}$ using the conjunction encoding above. The complete selection model is
\begin{align}
 \min\quad &T,\nonumber\\
 \text{subject to}\quad
 &z_{wh\mu c}\le K_{wh\mu c} &&\forall w,h,\mu,c,\nonumber\\
 &\sum_h A_hy_h+\sum_{h,\mu}A^B_{h\mu}v_{h\mu}\le A_{\max},\label{eq:finite-boundary}\\
 &L_w=\sum_{h,\mu,c}t_{wh\mu c}z_{wh\mu c} &&\forall w,\nonumber\\
 &T\ge L_w/L_w^{\mathrm{ref}} &&\forall w.\nonumber
\end{align}
Hardware and boundary mode are shared, while programs differ across scenarios. The area of an asynchronous engine is paid once even when several scenario programs use it. Its per-scenario effect is represented in $t$ and $K$. Latency entries must include startup, queue occupancy, synchronization, and the attainable overlap for the selected variant.

This finite model is an exact optimizer of its supplied table. Its purpose is to make the boundary decision and accounting explicit. Scaling to a large family replaces parts of the table with factorized resource and event constraints, and generates additional variants as needed. Doing so changes the computational representation while retaining the shared implementation identities. A boundary experiment would populate these coefficients from generated software and hardware; the empirical studies in Section~\ref{sec:evaluation} retain their existing instruction interface.

\section{Empirical Definitions and Source Records}\label{app:evidence}
The accompanying artifact stores the public source reports and comments as retrieved on October 9, 2026, alongside machine-readable data used in the figures. These are frozen report snapshots, not a substitute for executable experiment packages. The source reports identify the underlying branches and available reproduction material. No new accelerator simulation is required to reproduce the manuscript figures.

\textbf{Controlled study.}
The objectives in Figure~\ref{fig:controlled}b are the maximum normalized cycles across two prefill or two decode scenarios. In the order baseline, software-only, hardware-only, software-first, joint, and alternating, the prefill values are
$1$, $0.995747$, $0.908609$, $0.895792$, $0.891559$, and $0.891559$.
The decode values are
$1$, $0.816048$, $0.949019$, $0.813232$, $0.779915$, and $0.779915$.
The relative gain over software-first is $(r_{\mathrm{seq}}-r_{\mathrm{joint}})/r_{\mathrm{seq}}$.
The source is the controlled study~\cite{ventus-controlled}, with the snapshot retaining its update timestamp.

The study's acquisition domain contains 735 hardware profiles; seven configurations form the final evaluated selection menu. The report records 584 newly completed native evaluations and approximately 149.7 minutes of campaign wall time. Cache reuse and historical receipts increase the final data index to 1568 entries; that count is not the number of new simulations. The final menu's zero MIP gap describes its selection model. Representative-program aggregation, sampled decode steps, and ideal resource replication define the whole-network estimate.

\textbf{Tiling trajectory.}
Figure~\ref{fig:trajectory} uses the earlier accepted prefill objective 0.367 from the issue body, N3's rounded objective 0.342 from comment 6077807611, and N4's 0.3417 from comment 6077879211~\cite{ventus-tiling}. These are successive accepted designs, not equal-budget policy trials. The N4 ratios are $0.3417$, $0.3308$, $0.2953$, and $0.2719$; its estimated area is 1.0812\,mm$^2$. The figure preserves the precision available in each source entry.

\textbf{Model fidelity and transfer.}
The model study~\cite{ventus-model} uses complete Qwen instruction-timing runs and fixed-hardware transfer tests. Its software-only option remains the baseline software. The final configurations reported for transfer are the best designs among each task's specified controls, not necessarily the design labeled joint. GPT-2 small and Pythia-160M use FP16 software paths, whereas Qwen uses BF16. Their cycle ratios are normalized to the corresponding workload's own baseline, and are not pooled with the 14B calibrated study.
